\subsection{Superimposed integrals of positive functions}

For the rest of the section, absent express mention to the contrary, we shall denote by X a locally compact space, by \(\Lambda: t \mapsto \lambda_{t}\) a \(\mu\) -adequate mapping of T into \(\mathcal{M}_{+}(X)\) , and by \(\nu\) the integral of \(\Lambda\) .

PROPOSITION 3. --- Let \(f\) be a numerical function \(\geqslant 0\) defined on X. a) The following inequalities hold:

\[
\int^ {*} f (x) d \nu (x) \geqslant \int^ {\bullet} d \mu (t) \int^ {*} f (x) d \lambda_ {t} (x) \geqslant \int^ {\bullet} d \mu (t) \int^ {\bullet} f (x) d \lambda_ {t} (x). \tag {6}
\]

\begin{enumerate}
\def\labelenumi{\alph{enumi})}
\setcounter{enumi}{1}
\tightlist
\item
  If \(\Lambda\) is vaguely continuous, then
\end{enumerate}

\[
\int^ {*} f (x) d \nu (x) \geqslant \int^ {*} d \mu (t) \int^ {*} f (x) d \lambda_ {t} (x).\tag{7}
\]

\begin{enumerate}
\def\labelenumi{\alph{enumi})}
\setcounter{enumi}{2}
\tightlist
\item
  If \(\lambda_t^\bullet (1) <   + \infty\) locally \(\mu\) -almost everywhere, then
\end{enumerate}

\[
\int^ {\bullet} f (x) d \nu (x) \geqslant \int^ {\bullet} d \mu (t) \int^ {*} f (x) d \lambda_ {t} (x) = \int^ {\bullet} d \mu (t) \int^ {\bullet} f (x) d \lambda_ {t} (x). \tag {8}
\]

Let \(g\) be a lower semi-continuous function on \(X\) such that \(f \leqslant g\) . For every \(t \in T\) ,

\[
\int^ {*} f (x) d \lambda_ {t} (x) \leqslant \int^ {*} g (x) d \lambda_ {t} (x),
\]

therefore, by (4) and Prop. 4 of §1,

\[
\int^ {\bullet} d \mu (t) \int^ {*} f (x) d \lambda_ {t} (x) \leqslant \int^ {\bullet} d \mu (t) \int^ {*} g (x) d \lambda_ {t} (x) = \int^ {*} g (x) d \nu (x).
\]

The first of the inequalities (6) then follows from the definition of \(\int^{*}f(x)d\nu(x)\) (Ch. IV, §1, No.~3, Def. 3), and the second follows immediately from it. The inequality (7) is proved in an analogous way if \(\Lambda\) is vaguely continuous, using (5) instead of (4).

Let us pass to the proof of (8). The mapping \(t \mapsto \lambda_t^*(1)\) is measurable, and finite locally \(\mu\) -almost everywhere. The set \(\mathfrak{K}\) of compact subsets of \(T\) such the restriction of \(t \mapsto \lambda_t^*(1)\) to \(K\) is finite and continuous is therefore \(\mu\) -dense, and Prop. 4 of §2, No.~3 implies the existence of a summable family \((\mu_\alpha)_{\alpha \in A}\) of positive measures, with supports belonging to \(\mathfrak{K}\) , such that \(\mu = \sum_{\alpha \in A} \mu_\alpha\) . The mapping \(\Lambda\) is \(\mu_\alpha\) -adequate for every \(\alpha \in A\) ; set \(\nu_\alpha = \int \lambda_t d\mu_\alpha(t)\) . Prop. 1 shows that \(\nu = \sum_{\alpha \in A} \nu_\alpha\) , and the relation (4), applied to the measure \(\mu_\alpha\) and the function 1, shows that \(\nu_\alpha\) is a bounded measure (because \(\lambda_t^*(1)\) is bounded on \(\operatorname{Supp}(\mu_\alpha)\) ). Let us then write the formula (6) for the measure \(\mu_\alpha\) , replacing the symbol \(\int^*\) in the first member by \(\int^\bullet\) , which is legitimate by Prop. 7 of §1; then

\[
\int^ {\bullet} f (x) d \nu_ {\alpha} (x) \geqslant \int^ {\bullet} d \mu_ {\alpha} (t) \int^ {*} f (x) d \lambda_ {t} (x) = \int^ {\bullet} d \mu_ {\alpha} (t) \int^ {\bullet} f (x) d \lambda_ {t} (x)
\]

(the last equality due to the fact that \(\lambda_t\) is bounded locally almost everywhere, and Prop. 7 of § 1). The inequality in (8) is then obtained by summing on \(\alpha\) (§2, No.~2, Prop. 1).

If no hypothesis analogous to that of c) is made, the inequality (8) may fail (Exer. 2).

COROLLARY 1. --- Let \(f\) be a function \(\geqslant 0\) defined on \(X\) , and let \(H\) be the set of \(t \in T\) such that \(f\) is not \(\lambda_t\) -negligible.

\begin{enumerate}
\def\labelenumi{\alph{enumi})}
\item
  If \(f \nu\) -negligible, then \(H\) is locally \(\mu\) -negligible.
\item
  If \(f\) is \(\nu\) -negligible and \(\Lambda\) is vaguely continuous, then \(H\) is \(\mu\) -negligible.
\item
  If \(f\) is locally \(\nu\) -negligible and \(\lambda_t^\bullet(1) < +\infty\) locally \(\mu\) -almost everywhere, then \(H\) is locally \(\mu\) -negligible.
\end{enumerate}

COROLLARY 2. --- Let f be a function \(\geqslant0\) defined on X, \(\nu\) -measurable and \(\nu\) -moderated. The set of \(t\in T\) such that f is not \(\lambda_{t}\) -moderated is then locally \(\mu\) -negligible (and even \(\mu\) -negligible if \(\Lambda\) is vaguely continuous).

For, \(f\) is the sum of a sequence of functions \(f_{n} \geqslant 0\) such that \(f_{n}\) is zero outside a compact set \(K_{n}\) for \(n \geqslant 1\) , and \(f_{0}\) is \(\nu\) -negligible (§1, No.~2, Prop. 6); \(f_{0}\) is then \(\lambda_{t}\) -negligible except for \(t\) forming a set that is locally \(\mu\) -negligible (and even \(\mu\) -negligible, if \(\Lambda\) is vaguely continuous) by Cor. 1, and the statement then follows at once.

PROPOSITION 4. --- Let f be a \(\nu\) -measurable function defined on X, with values in a topological space G, and let M be the set of \(t \in T\) such that f is not \(\lambda_{t}\) -measurable.

\begin{enumerate}
\def\labelenumi{\alph{enumi})}
\item
  Suppose that \(f\) is constant on the complement of a \(\nu\) -moderated subset of \(X\) ; then \(M\) is locally \(\mu\) -negligible.
\item
  Suppose that \(f\) is constant on the complement of a \(\nu\) -moderated subset of \(X\) , and that \(\Lambda\) is vaguely continuous; then \(M\) is \(\mu\) -negligible.
\item
  Suppose that \(\lambda_t^\bullet(1) < +\infty\) locally \(\mu\) -almost everywhere; then M is locally \(\mu\) -negligible.
\end{enumerate}

Let us first prove a) (resp. b)). Since every \(\nu\) -integrable set is contained in a \(\nu\) -integrable open set, the function f is constant on the complement B of a countable union of \(\nu\) -integrable open sets. There exists a partition of X - B formed by a \(\nu\) -negligible set N and a sequence \((\mathrm{K}_{n})\) of compact sets such that the restriction of f to each \(K_{n}\) is continuous. Let S be the set of \(t \in T\) such that N is not \(\lambda_{t}\) -negligible: S is locally \(\mu\) -negligible (resp. \(\mu\) -negligible) by Cor. 1 of Prop. 3. The sets \(K_{n}, B, N\) are measurable for every measure on X, and the restriction of f to each of them is \(\lambda_{t}\) -measurable for every \(t \notin S\) . The function f is therefore \(\lambda_{t}\) -measurable for every \(t \notin S\) (Ch. IV, §5, No.~10, Prop. 16).

To establish c), let us take up again the notations in the proof of Prop. 3; since f is \(\nu\) -measurable, it is measurable for each of the measures \(\nu_{\alpha} \leqslant \nu\) . Now, these measures are bounded, hence moderated, and it follows from a) that M is locally \(\mu_{\alpha}\) -negligible for every \(\alpha \in A\) . This implies that M is locally \(\mu\) -negligible (§2, No.~2, Cor. 2 of Prop. 1).

PROPOSITION 5. --- Let f be a \(\nu\) -measurable positive numerical function defined on X, and let N be the set of \(t \in T\) such that f is not both \(\lambda_{t}\) -measurable and \(\lambda_{t}\) -moderated.

\begin{enumerate}
\def\labelenumi{\alph{enumi})}
\tightlist
\item
  Suppose that \(f\) is \(\nu\) -moderated. The set \(\mathbf{N}\) is then locally \(\mu\) -negligible, the function \(t \mapsto \int^{\bullet} f(x) d\lambda_t(x)\) is \(\mu\) -measurable, and
\end{enumerate}

\[
\int^ {\bullet} f (x) d \nu (x) = \int^ {\bullet} d \mu (t) \int^ {\bullet} f (x) d \lambda_ {t} (x).\tag{9}
\]

\begin{enumerate}
\def\labelenumi{\alph{enumi})}
\setcounter{enumi}{1}
\tightlist
\item
  Suppose that \(f\) is \(\nu\) -moderated, and that \(\Lambda\) is vaguely continuous. The set \(N\) is then \(\mu\) -negligible, the function \(t \mapsto \int^{*} f(x) d\lambda_{t}(x)\) is \(\mu\) -measurable and \(\mu\) -moderated, and
\end{enumerate}

\[
\int^ {*} f (x) d \nu (x) = \int^ {*} d \mu (t) \int^ {*} f (x) d \lambda_ {t} (x).\tag{10}
\]

\begin{enumerate}
\def\labelenumi{\alph{enumi})}
\setcounter{enumi}{2}
\tightlist
\item
  Suppose that \(\lambda_t^\bullet(1) < +\infty\) locally \(\mu\) -almost everywhere. The set N is then locally \(\mu\) -negligible, the function \(t \mapsto \int^\bullet f(x) d\lambda_t(x)\) is \(\mu\) -measurable, and (9) holds.
\end{enumerate}

Let us first prove a) (resp. b)), assuming that \(f\) is \(\nu\) -moderated. The assertions concerning the set \(N\) have already been established (Prop. 4, and Cor. 2 of Prop. 3). By Prop. 6 of §1, No.~2, we may limit ourselves to proving a) (resp. b)) in each of the following special cases:

\begin{enumerate}
\def\labelenumi{\arabic{enumi})}
\item
  The function \(f\) is \(\nu\) -negligible.
\item
  There exists a compact set \(\mathbf{K}\) such that \(f\) is zero outside \(\mathbf{K}\) and the restriction of \(f\) to \(\mathbf{K}\) is continuous.
\end{enumerate}

The special case 1) has already been treated (Cor. 1 of Prop. 3). To treat the second, we denote by G a \(\nu\) -integrable open set containing K, by M a constant upper bound for f, by h the lower semi-continuous function \(M\varphi_{G}\) , and by g the function h-f.~Since the function f is upper semi-continuous on X, g is lower semi-continuous and positive. Moreover, f,g,h are \(\nu\) -integrable.

Let us then apply formula (4) (resp. (5)) to the lower semi-continuous functions h and g. By subtraction, we see that the function

\[
t \mapsto \int^ {\bullet} f (x) d \lambda_ {t} (x) \quad (\text { resp. } \int^ {*} f (x) d \lambda_ {t} (x))
\]

is \(\mu\) -measurable and that the formula (9) (resp. (10)) holds. Finally, under the hypothesis of b), the function \(t \mapsto \int^{*} f(x) d\lambda_{t}(x)\) has finite upper integral, hence is indeed \(\mu\) -moderated.

To prove c), let us take up again the measures \(\mu_{\alpha}\) and \(\nu_{\alpha}\) of the proof of Prop. 3; since f is \(\nu_{\alpha}\) -measurable and \(\nu_{\alpha}\) -moderated, the assertion a) implies that \(t \mapsto \int^{\bullet} f(x) d\lambda_{t}(x)\) is \(\mu_{\alpha}\) -measurable and that

\[
\int^ {\bullet} f (x) d \nu_ {\alpha} (x) = \int^ {\bullet} d \mu_ {\alpha} (t) \int^ {\bullet} f (x) d \lambda_ {t} (x).
\]

It remains only to sum on \(\alpha\) , applying Props. 1 and 2 of §2, No.~2.

If \(f\) is not assumed to be \(\nu\) -moderated, and if one does not make the assumption in c), then the relation (9) may not hold (Exer. 3).

COROLLARY. --- Let f be a function defined on X, with values in a Banach space F or in \(\overline{R}\) , that is \(\nu\) -measurable and \(\nu\) -moderated. For f to be \(\nu\) -integrable, it is necessary and sufficient that

\[
\int^ {\bullet} d \mu (t) \int^ {\bullet} | \mathbf {f} (x) | d \lambda_ {t} (x) <   + \infty .
\]

This follows at once from Prop. 5 and the criterion for integrability (Ch. IV, §5, No.~6, Th. 5).

\subsection{Superimposed integrals of functions with values in a Banach space}

THEOREM 1. --- Let f be a function with values in a Banach space F or in \(\overline{R}\) , and let H be the set of \(t \in T\) for which f is not \(\lambda_{t}\) -integrable.

\begin{enumerate}
\def\labelenumi{\alph{enumi})}
\tightlist
\item
  If \(\mathbf{f}\) is \(\nu\) -integrable, then \(\mathrm{H}\) is locally \(\mu\) -negligible, the function \(t \mapsto \int \mathbf{f}(x)d\lambda_t(x)\) (defined for \(t \notin \mathrm{H}\) ) is essentially \(\mu\) -integrable, and
\end{enumerate}

\[
\int \mathbf {f} (x) d \nu (x) = \int d \mu (t) \int \mathbf {f} (x) d \lambda_ {t} (x).\tag{11}
\]

\begin{enumerate}
\def\labelenumi{\alph{enumi})}
\setcounter{enumi}{1}
\item
  If \(\mathbf{f}\) is \(\nu\) -integrable and \(\Lambda\) is vaguely continuous, then \(\mathbf{H}\) is moreover \(\mu\) -negligible and the function \(t \mapsto \int \mathbf{f}(x) d\lambda_t(x)\) (defined for \(t \notin \mathbf{H}\) ) is \(\mu\) -integrable.
\item
  If \(\lambda_t^\bullet(1) < +\infty\) locally \(\mu\) -almost everywhere, then the conclusions of a) remain true for f an essentially \(\nu\) -integrable function.
\end{enumerate}

We are first going to establish a) (resp. b)). This statement is true when f is a positive numerical function (Prop. 5); if f is an integrable function with values in \(\overline{R}\) , this result may be applied to the positive functions \(f^{+}\) and \(f^{-}\) , and therefore extends at once to f by subtraction. It remains to treat the case of functions with values in F. Let H be the subspace of \(\mathcal{L}_{\mathrm{F}}^{1}(\nu)\) formed by the linear combinations, with coefficients in F, of the functions in \(\mathcal{H}(\mathrm{X})\) ; the result pertaining to real functions implies at once the validity of the statement for the elements of H. Now, H is dense in \(\mathcal{L}_{\mathrm{F}}^{1}(\nu)\) ; therefore, for every \(\mathbf{f} \in \mathcal{L}_{\mathrm{F}}^{1}(\nu)\) , there exists a sequence \((\mathbf{f}_{n})\) of elements of H that has the following properties:

\begin{enumerate}
\def\labelenumi{\arabic{enumi})}
\item
  the sequence \((\mathbf{f}_n)\) converges to \(\mathbf{f}\) in mean in \(\mathcal{L}_{\mathrm{F}}^{1}(\nu)\) , and \(\nu\) -almost everywhere;
\item
  the function \(g = |\mathbf{f}_0| + \sum_{n\in \mathbf{N}}|\mathbf{f}_{n + 1} - \mathbf{f}_n|\) is such that \(\nu^{*}(g) < +\infty\) (Ch. IV, §3, No.~4, Th. 3).
\end{enumerate}

Let \(N_{1}\) be the set of \(t \in T\) such that \(\lambda_{t}^{*}(g) = +\infty\) ; \(N_{1}\) is locally \(\mu\) -negligible (resp. \(\mu\) -negligible) by formula (6) (resp. (7)). For \(t \notin N_{1}\) , the \(f_{n}\) belong to \(\mathcal{L}_{F}^{1}(\lambda_{t})\) , the sequence \((f_{n})\) converges \(\lambda_{t}\) -almost everywhere, as well as for the topology of convergence in mean in \(\mathcal{L}_{F}^{1}(\lambda_{t})\) (Ch. IV, §3, No.~3, Prop. 6). Let M be set of \(x \in X\) such that \(f_{n}(x)\) does not converge to \(f(x)\) : since M is \(\nu\) -negligible, the set \(N_{2}\) of \(t \in T\) such that M is not \(\lambda_{t}\) -negligible is locally \(\mu\) -negligible (resp. \(\mu\) -negligible) by Cor. 1 of Prop. 3.

Suppose that \(t\) does not belong to \(\mathrm{N}_1 \cup \mathrm{N}_2\) ; the sequence \((\mathbf{f}_n)\) converges in mean in \(\mathcal{L}_{\mathrm{F}}^{1}(\lambda_t)\) , and converges \(\lambda_t\) -almost everywhere to \(\mathbf{f}\) . Therefore \(\mathbf{f} \in \mathcal{L}_{\mathrm{F}}^{1}(\lambda_t)\) and \(\int \mathbf{f} d\lambda_t = \lim_{n \to \infty} \int \mathbf{f}_n d\lambda_t\) (Ch. IV, §4, No.~1). The set H of the statement is thus contained in \(N_{1} \cup N_{2}\) ; it is therefore locally \(\mu\) -negligible (resp. \(\mu\) -negligible). On the other hand, the function \(t \mapsto \int f d\lambda_{t}\) is equal locally \(\mu\) -almost everywhere to the limit of a sequence of \(\mu\) -measurable functions; it is therefore \(\mu\) -measurable. Finally, for every \(t \notin N_{1} \cup N_{2}\) and every n, we have

\[
\left| \int \mathbf {f} _ {n} (x) d \lambda_ {t} (x) \right| \leqslant \int^ {*} g (x) d \lambda_ {t} (x)
\]

by virtue of the inequality \(|\mathbf{f}_n| \leqslant g\) and Prop. 2 of Ch. IV, §4, No.~2. Now, the function \(t \mapsto \int^{*} g(x) d\lambda_t(x)\) is essentially \(\mu\) -integrable (resp. \(\mu\) -integrable) by Prop. 5. We may therefore apply Lebesgue's theorem, which yields

\[
\int d \mu (t) \int \mathbf {f} (x) d \lambda_ {t} (x) = \lim _ {n \rightarrow \infty} \int d \mu (t) \int \mathbf {f} _ {n} (x) d \lambda_ {t} (x) = \lim _ {n \rightarrow \infty} \int \mathbf {f} _ {n} (x) d \nu (x).
\]

Since \(\int \mathbf{f}_n(x) d\nu(x)\) tends to \(\int \mathbf{f}(x) d\nu(x)\) as \(n\) tends to \(\infty\) , by the hypotheses made on the sequence \((\mathbf{f}_n)\) , the relation (11) follows and we have proved a) (resp. b)).

Now suppose that \(\lambda_t^\bullet(1) < +\infty\) locally \(\mu\) -almost everywhere, and that \(\mathbf{g}\) is an essentially \(\nu\) -integrable function. Let \(\mathbf{f}\) be a \(\nu\) -integrable function such that \(\mathbf{g} = \mathbf{f}\) locally \(\nu\) -almost everywhere (§1, No.~3). Then \(\mathbf{g} = \mathbf{f}\) almost everywhere for \(\lambda_t\) , except for \(t\) forming a locally \(\mu\) -negligible set \(\mathrm{P}\) (Cor. 1 c) of Prop. 3). Therefore \(\int \mathbf{g} d\lambda_t = \int \mathbf{f} d\lambda_t\) for all \(t \notin \mathrm{P} \cup \mathrm{H}\) , and this completes the proof.

Remark. --- Let \(\Lambda: t \mapsto \lambda_{t}\) be a \(\mu\) -adequate mapping of T into \(\mathcal{M}_{+}(X)\) . If a mapping \(\Lambda': t \mapsto \lambda_{t}'\) of T into \(\mathcal{M}_{+}(X)\) is equal to \(\Lambda\) locally \(\mu\) -almost everywhere, it follows at once from the definitions that \(\Lambda'\) is also \(\mu\) -adequate, and that \(\Lambda\) and \(\Lambda'\) have the same integral. If now \(H: t \mapsto \eta_{t}\) is a function with values in \(\mathcal{M}_{+}(X)\) , defined locally \(\mu\) -almost everywhere, we shall again say that H is \(\mu\) -adequate if it is equal locally \(\mu\) -almost everywhere to a mapping \(\Lambda: t \mapsto \lambda_{t}\) that is everywhere defined and \(\mu\) -adequate. We then set \(\int \eta_{t} d\mu(t) = \int \lambda_{t} d\mu(t)\) , a definition that does not depend on the function \(\Lambda\) utilized. We leave to the reader the task of verifying that the propositions proved in the preceding Nos. extend to \(\mu\) -adequate functions defined locally \(\mu\) -almost everywhere.

\subsection{Universally measurable functions}

DEFINITION 2. --- A mapping f of T into a topological space F is said to be universally measurable if it is \(\mu\) -measurable for every positive measure \(\mu\) on T.

The subsets of T whose characteristic function is universally measurable are called universally measurable sets. They form a tribe on T (Ch. IV, §5, No.~4, Cor. 2 of Th. 2) that contains the Borel sets (same ref., Cor. 3), and the Souslin sets if T is metrizable (Ch. IV, §5, No.~1, Cor. 2 of Prop. 3). For a mapping f of T into a topological space F, metrizable and separable, to be universally measurable, it is necessary and sufficient that the inverse image under f of every closed ball in F be a universally measurable subset of T (Ch. IV, §5, No.~5, Th. 4).

PROPOSITION 6. --- For a mapping f of T into a topological space F to be universally measurable, it is necessary and sufficient that f be measurable for every positive measure on T with compact support.

The condition is obviously necessary; on the other hand it is sufficient, because every positive measure \(\mu\) is the sum of a family of measures with compact support ( \(\S2\) , No.~3, Prop. 4): the statement then follows from Prop. 2 of \(\S2\) , No.~2.

PROPOSITION 7. --- Let \(\mu\) be a positive measure on T, and let \(f\) be a \(\mu\) -measurable mapping of T into a topological space F. Then, there exists a universally measurable mapping \(f'\) of T into F such that \(f = f'\) locally \(\mu\) -almost everywhere.

Let \(\mathfrak{K}\) be the set of compact sets in T such that the restriction of \(f\) to K is continuous; since \(\mathfrak{K}\) is \(\mu\) -dense (Ch. IV, §5, No.~10, Prop. 15), there exists a locally countable family \((\mathrm{K}_i)_{i\in I}\) of pairwise disjoint elements of \(\mathfrak{K}\) such that the set \(\mathrm{N} = \mathrm{T} - \bigcup_{i\in I}\mathrm{K}_i\) is locally \(\mu\) -negligible (Ch. IV, §5, No.~9,

Prop. 14). Let \(x\) be an element of \(\mathbf{F}\) ; set

\[
f ^ {\prime} (t) = f (t) \quad \text { if } t \in \bigcup_ {i \in I} K _ {i},
\]

\[
f ^ {\prime} (t) = x \quad \text { if } t \in \mathbb {N}.
\]

The functions \(f\) and \(f'\) are equal locally \(\mu\) -almost everywhere. On the other hand, \(\mathrm{N} \cap \mathrm{K}\) is a Borel subset of \(\mathrm{K}\) for every compact set \(\mathrm{K}\) in \(\mathrm{T}\) , since the family \((\mathrm{K}_i)\) is locally countable. It follows that \(\mathrm{N}\) is a universally measurable set, and that \(f'\) is a universally measurable function (Ch. IV, §5, No.~10, Prop. 16).

\subsection{Diffusions}

DEFINITION 3. --- Let X be a locally compact space, and let \(\Lambda : t \mapsto \lambda_t\) be a mapping of T into \(\mathcal{M}_{+}(X)\) . The mapping \(\Lambda\) is said to be a diffusion of T in X if \(\Lambda\) is adequate for every positive measure on T with compact support. The diffusion \(\Lambda\) is said to be bounded if all of the measures \(\lambda_{t}\) are bounded and \(\sup_{t\in T}\|\lambda_{t}\|<+\infty\) ; this quantity is then called the norm of \(\Lambda\) and is denoted \(\|\Lambda\|\) .

The following proposition merely translates the definition:

PROPOSITION 8. --- For a mapping \(\Lambda: t \mapsto \lambda_t\) of T into \(\mathcal{M}_+(\mathrm{X})\) to be a diffusion, it is necessary and sufficient that the following conditions be satisfied:

\begin{enumerate}
\def\labelenumi{\arabic{enumi})}
\item
  For every lower semi-continuous function \(f \geqslant 0\) defined on \(X\) , the function \(t \mapsto \lambda_t^\bullet(f)\) is universally measurable on \(T\) .
\item
  For every function \(g \in \mathcal{K}_{+}(X)\) , the function \(t \mapsto \lambda_t(g)\) is locally bounded in \(T\) .
\item
  For every lower semi-continuous function \(f \geqslant 0\) defined on X and for every positive measure \(\mu\) on T with compact support, the following relation holds, where \(\nu\) denotes \(\int \lambda_{t} d\mu(t)\) :
\end{enumerate}

\[
\int^ {\bullet} f (x) d \nu (x) = \int^ {\bullet} d \mu (t) \int^ {\bullet} f (x) d \lambda_ {t} (x).\tag{12}
\]

Suppose that \(\Lambda\) is a diffusion. The condition 1) is then satisfied by the definition of adequate mappings (No.~1, Def. 1) and Prop. 6; the condition 3) is satisfied by formula (4), since \(\Lambda\) is \(\mu\) -adequate. Let \(g \in \mathcal{K}_{+}(X)\) and let \(u\) be the function \(t \mapsto \lambda_t(g)\) (universally measurable, by 1)); suppose that \(u\) is not locally bounded. There would then exist a compact set \(K\) such that \(u\) is not bounded on \(K\) , hence there would exist a sequence \((t_n)\) of elements of \(K\) such that \(u(t_n) \geqslant n^2\) for all \(n \geqslant 1\) ; then \(u\) is not integrable for the measure \(\mu = \sum_{n \geqslant 1} \frac{1}{n^2} \varepsilon_{t_n}\) with compact support, contrary to the hypothesis on \(\Lambda\) , which implies that \(t \mapsto \lambda_t(g)\) is integrable for every positive measure with compact support. The above three conditions are thus necessary. Conversely, the conditions 1) and 2) imply that \(\Lambda\) is scalarly essentially \(\mu\) -integrable for every measure \(\mu\) with compact support. Since every measure \(\mu' \geqslant 0\) that is bounded above by a measure \(\mu\) with compact support also has compact support, the conditions 1) and 3) express that \(\Lambda\) is \(\mu\) -adequate for every positive measure with compact support, which is indeed the sought-for result.

PROPOSITION 9. --- Let \(\Lambda : t \mapsto \lambda_t\) be a mapping of T into \(\mathcal{M}_+(\mathrm{X})\) , such that the function \(t \mapsto \lambda_t(g)\) is universally measurable and locally bounded in T for every \(g \in \mathcal{K}_+(\mathrm{X})\) . One can affirm that \(\Lambda\) is a diffusion in each of the following cases:

\begin{enumerate}
\def\labelenumi{\alph{enumi})}
\item
  the topology of X admits a countable base;
\item
  \(\Lambda\) is universally measurable for the vague topology.
\end{enumerate}

For, let \(\mu\) be a positive measure on T with compact support; the mapping \(\Lambda\) is scalarly essentially \(\mu\) -integrable, hence \(\mu\) -adequate if either a) or b) is satisfied (No.~1, Prop. 2).

For the rest of this section, we shall adopt the following notations: we will denote by \(\langle\eta,h\rangle\) the upper essential integral, for a positive measure \(\eta\) , of a positive \(\eta\) -measurable function h. The mapping \(\Lambda:t\mapsto\lambda_{t}\) will be a diffusion of T in X. If f is a positive universally measurable function defined on X, we shall denote by \(\Lambda f\) the mapping \(t\mapsto\lambda_{t}^{\bullet}(f)\) . If \(\mu\) is a positive measure on T such that \(\Lambda\) is scalarly essentially \(\mu\) -integrable, we shall denote by \(\mu\Lambda\) the measure \(\int\lambda_{t}d\mu(t)\) . The definition of the integral then takes the form

\[
\langle \mu \Lambda , f \rangle = \langle \mu , \Lambda f \rangle \quad \text { for } f \in \mathscr {K} _ {+} (\mathrm{X}).
\]

We shall say that a positive measure \(\mu\) on T belongs to the domain of \(\Lambda\) if \(\Lambda\) is \(\mu\) -adequate: this amounts to saying (in view of Prop. 8) that \(\Lambda\) is scalarly essentially \(\mu\) -integrable and \(\langle\mu'\Lambda,f\rangle=\langle\mu',\Lambda f\rangle\) for every positive measure \(\mu'\leqslant\mu\) and every lower semi-continuous positive function f.

PROPOSITION 10. --- Let f, g be two positive universally measurable functions on X, let a be a number \(\geqslant 0\) , and let \(\mu\) and \(\nu\) be two positive measures on T. Then:

\begin{enumerate}
\def\labelenumi{\alph{enumi})}
\item
  \(\Lambda (f + g) = \Lambda f + \Lambda g, \Lambda (af) = a\Lambda f.\)
\item
  If \(\mu\) and \(\nu\) belong to the domain of \(\Lambda\) , then so do \(\mu + \nu\) and \(a\mu\) , and one has \((\mu + \nu)\Lambda = \mu \Lambda + \nu \Lambda\) , \((a\mu)\Lambda = a(\mu \Lambda)\) .
\end{enumerate}

The only non-obvious point is that \(\mu + \nu\) belongs to the domain of \(\Lambda\) , which is treated by observing that every positive measure bounded above by \(\mu + \nu\) is of the form \(\mu' + \nu'\) , where \(\mu' \leqslant \mu\) , \(\nu' \leqslant \nu\) (the `decomposition lemma', Ch. II, §1, No.~1). See also the next proposition.

PROPOSITION 11. --- For a positive measure \(\mu\) on \(T\) to belong to the domain of \(\Lambda\) , it is necessary and sufficient that \(\Lambda\) be scalarly essentially \(\mu\) -integrable.

This condition is obviously necessary. Conversely, suppose it is satisfied, and let f be a lower semi-continuous positive function defined on X. The function \(\Lambda f\) is universally measurable, hence \(\mu\) -measurable. We are going to prove that \(\langle\mu,\Lambda f\rangle=\langle\mu\Lambda,f\rangle\) ; since this equality will also be valid for every positive measure \(\mu'\leqslant\mu\) , because \(\Lambda\) is also scalarly essentially \(\mu'\) -integrable, it will follow that \(\Lambda\) is \(\mu\) -adequate.

Let \((\mu_i)_{i\in I}\) be a summable family of positive measures with compact support, such that \(\mu = \sum_{i\in I}\mu_i\) (§2, No.~3, Prop. 4); the family of measures \(\mu_{i}\Lambda\) is then summable, and \(\mu \Lambda = \sum_{i\in I}\mu_{i}\Lambda\) (No.~1, Cor. of Prop. 1). Consequently \(\langle \mu \Lambda ,f\rangle = \sum_{i\in I}\langle \mu_i\Lambda ,f\rangle\) (§2, No.~2, Prop. 1); but \(\Lambda\) is \(\mu_{i}\) -adequate, thus \(\langle \mu_i\Lambda ,f\rangle = \langle \mu_i,\Lambda f\rangle\) . Again applying Prop. 1 of §2, we obtain the sought-for equality:

\[
\langle \mu \Lambda , f \rangle = \sum_ {i \in I} \langle \mu_ {i} \Lambda , f \rangle = \sum_ {i \in I} \langle \mu_ {i}, \Lambda f \rangle = \langle \mu , \Lambda f \rangle .
\]

COROLLARY 1. --- If \(\Lambda\) is a bounded diffusion, then every bounded positive measure \(\mu\) belongs to the domain of \(\Lambda\) , and \(\|\mu\Lambda\| \leqslant \|\mu\|\|\Lambda\|\) .

COROLLARY 2. --- Suppose that \(\mu\) is the sum of a summable family \((\mu_{\alpha})_{\alpha \in \mathbb{A}}\) of positive measures belonging to the domain of \(\Lambda\) . For \(\mu\) to belong to the domain of \(\Lambda\) , it is necessary and sufficient that the family of measures \(\mu_{\alpha}\Lambda\) be summable, in which case \(\mu \Lambda = \sum \mu_{\alpha}\Lambda\) .

It suffices to apply the Corollary of Prop. 1 of No.~1.

Proposition 5, expressed in the language of diffusions, takes the following form:

PROPOSITION 12. --- Let \(\mu\) be a positive measure on T that belongs to the domain of \(\Lambda\) , and let \(f\) be a universally measurable function \(\geqslant 0\) defined on X. If \(f\) is moderated for the measure \(\mu\Lambda\) , or if the measures \(\lambda_t\) are bounded, then the function \(\Lambda f\) is \(\mu\) -measurable and

\[
\langle \mu \Lambda , f \rangle = \langle \mu , \Lambda f \rangle .\tag{13}
\]

COROLLARY. --- If X is countable at infinity, or if the measures \(\lambda_{t}\) are bounded, then the function \(\Lambda f\) is universally measurable on T for every universally measurable function \(f \geqslant 0\) defined on X, and (13) holds.

\subsection{Composition of bounded diffusions}

PROPOSITION 13. --- Let T,X,Y be three locally compact spaces, \(\Lambda : t \mapsto \lambda_t\) a bounded diffusion of T in X, and H: \(x \mapsto \eta_x\) a bounded diffusion of X in Y. The mapping \(t \mapsto \lambda_t\mathrm{H}\) is then a bounded diffusion of T in Y, which is denoted by \(\Lambda\mathrm{H}\) , and

\[
\| \Lambda H \| \leqslant \| \Lambda \| \| H \|.\tag{14}
\]

Let \(f\) be a universally measurable function \(\geqslant 0\) defined on \(Y\) , and \(\mu\) a measure on \(T\) . Suppose that \(\mu\) belongs to the domain of \(\Lambda\) , and that \(\mu \Lambda\) belongs to the domain of H; then \(\mu\) belongs to the domain of \(\Lambda \mathrm{H}\) , and

\[
\begin{array}{l} \langle \mu (\Lambda \mathrm{H}), f \rangle = \langle \mu \Lambda , \mathrm{H} f \rangle = \langle \mu , \Lambda \mathrm{H} f \rangle ; \\ (\mu \Lambda) \mathrm{H} = \mu (\Lambda \mathrm{H}); \quad \Lambda (\mathrm{H} f) = (\Lambda \mathrm{H}) f. \end{array}\tag{15}
\]

Set \(\gamma_{t} = \lambda_{t}H\) ; we shall denote by \(\Gamma\) the mapping \(\Lambda H\) of T into \(\mathcal{M}_{+}(Y)\) , and by \(\Gamma f\) the function \(t \mapsto \langle \gamma_{t}, f \rangle\) (an abuse of notation, since we do not yet know whether \(\Gamma\) is a diffusion). Then \(\langle \gamma_{t}, f \rangle = \langle \lambda_{t}H, f \rangle = \langle \lambda_{t}, Hf \rangle\) by (13); since the function Hf is positive and universally measurable on X (Cor. of Prop. 12), it follows first of all that \(\Gamma f = \Lambda(Hf)\) , and then that \(\Gamma f\) is universally measurable on T (same reference). It is clear that all of the measures \(\gamma_{t}\) have total mass at most equal to \(\|\Lambda\| \|H\|\) . Consequently \(\Gamma g\) is universally measurable and bounded for every function \(g \in \mathcal{K}_{+}(Y)\) ; \(\Gamma\) is therefore scalarly essentially integrable for every bounded measure on T, and in particular for every measure with compact support. More generally, if \(\mu\) is a measure in the domain of \(\Lambda\) , such that \(\mu\Lambda\) belongs to the domain of H, then, for \(g \in \mathcal{K}_{+}(Y)\) ,

\[
\langle \mu , \Gamma g \rangle = \langle \mu , \Lambda (\mathrm{H} g) \rangle = \langle \mu \Lambda , \mathrm{H} g \rangle = \langle (\mu \Lambda) \mathrm{H}, g \rangle .
\]

Since the last quantity is finite, we see that \(\Gamma\) is scalarly essentially \(\mu\) -integrable. Let us denote by \(\mu\Gamma\) the integral \(\int\gamma_{t}d\mu(t)\) (an abuse of notation, since we do not yet know whether \(\Gamma\) is a diffusion). The preceding relations may then be written

\[
\langle \mu \Gamma , g \rangle = \langle (\mu \Lambda) \mathrm{H}, g \rangle ,
\]

or also \(\mu \Gamma = (\mu \Lambda)\mathrm{H}\) since \(g\) is arbitrary in \(\mathcal{K}_+(\mathrm{Y})\) .

Consider anew the universally measurable function \(f \geqslant 0\) . We have

\[
\langle \mu \Gamma , f \rangle = \langle (\mu \Lambda) \mathrm{H}, f \rangle = \langle \mu \Lambda , \mathrm{H} f \rangle = \langle \mu , \Lambda (\mathrm{H} f) \rangle = \langle \mu , \Gamma f \rangle .
\]

When f is lower semi-continuous and \(\mu\) runs over the set of positive measures with compact support, these relations express that \(\Gamma\) is a diffusion of T in Y. The assertion then does no more than make explicit the relations obtained in the course of the above proof.

DEFINITION 4. --- The notations being those of Proposition 13, the diffusion \(\Lambda H\) is called the composed diffusion (or the composition) of the bounded diffusions \(H\) and \(\Lambda\) .

Let \(X_{1}, X_{2}, X_{3}, X_{4}\) be four locally compact spaces, and \(\Lambda_{1}, \Lambda_{2}, \Lambda_{3}\) three bounded diffusions, of \(X_{1}\) in \(X_{2}\) , \(X_{2}\) in \(X_{3}\) , \(X_{3}\) in \(X_{4}\) , respectively. It follows at once from Prop. 13 that

\[
(\Lambda_ {1} \Lambda_ {2}) \Lambda_ {3} = \Lambda_ {1} (\Lambda_ {2} \Lambda_ {3}).
\]

We will therefore use these notations without parentheses for the composition of diffusions.

Example. --- Let u be a universally measurable mapping of T into X, and v a universally measurable mapping of X into Y; by Prop. 2 b), one defines diffusions \(\Lambda\) and H by the formulas

\[
\lambda_ {t} = \varepsilon_ {u (t)}, \eta_ {x} = \varepsilon_ {v (x)};
\]

the diffusion \(\Gamma = \Lambda H\) is then given by

\[
\gamma_ {t} = \varepsilon_ {(v \circ u) (t)}.
\]

One is therefore careful to note that the order of composition of diffusions is the opposite of the usual order of the composition of functions.

\section{INTEGRATION OF POSITIVE POINT MEASURES}

\subsection{Families of point measures}

Let X and T be two locally compact spaces, \(\pi\) a mapping of T into X, and \(g\) a finite numerical function \(\geqslant 0\) defined on T; these two functions define a mapping \(t \mapsto \lambda_t = g(t)\varepsilon_{\pi(t)}\) of T into the space \(\mathcal{M}(X)\) of measures on X, such that for every \(t \in T\) , \(\lambda_t\) is either a point measure (Ch. III, §2, No.~4) or is equal to 0. If \(f\) is a numerical function \(\geqslant 0\) defined on X, then \(\int^{*} f(x) d\lambda_t(x) = \int^{\bullet} f(x) d\lambda_t(x) = f(\pi(t))g(t)\) (recall our convention of taking this product to be 0 when \(g(t) = 0\) and \(f(\pi(t)) = +\infty\) ). Every function (with values in a topological space) defined on X is \(\lambda_t\) -measurable for every \(t \in T\) . Every mapping \(\mathbf{f}\) of X into a Banach space F is \(\lambda_t\) -integrable for all \(t \in T\) , and \(\int \mathbf{f}(x) d\lambda_t(x) = \mathbf{f}(\pi(t))g(t)\) . Finally, if \(f\) is an arbitrary numerical function defined on X, for \(f\) to be \(\lambda_t\) -integrable it is necessary and sufficient that \(f(\pi(t))g(t)\) be finite, in which case \(\int f(x) d\lambda_t(x) = f(\pi(t))g(t)\) .

DEFINITION 1. --- Let \(\mu\) be a positive measure on T. The pair \((\pi, g)\) is said to be \(\mu\) -adapted if the following conditions are satisfied:

\(1^{\circ}\) The functions \(\pi\) and \(g\) are \(\mu\) -measurable.

\(2^{\circ}\) For every function \(f \in \mathcal{K}(X)\) , the mapping \(t \mapsto f(\pi(t))g(t)\) is essentially \(\mu\) -integrable.

PROPOSITION 1. --- If the pair \((\pi, g)\) is \(\mu\) -adapted, then the mapping \(\Lambda: t \mapsto \lambda_t = g(t) \varepsilon_{\pi(t)}\) of T into \(\mathcal{M}_+(\mathrm{X})\) is scalarly essentially \(\mu\) -integrable, vaguely \(\mu\) -measurable and \(\mu\) -adequate. Conversely, if \(\Lambda\) is scalarly essentially \(\mu\) -integrable and vaguely \(\mu\) -measurable, then the function g is \(\mu\) -measurable and the restriction of \(\pi\) to the set S of \(t \in T\) such that \(g(t) \neq 0\) is \(\mu\) -measurable.

Suppose that the pair \((\pi,g)\) is \(\mu\) -adapted; for every function \(f \in \mathcal{K}(X)\) , the function \(t \mapsto \langle f, \lambda_t \rangle = f(\pi(t))g(t)\) is then essentially \(\mu\) -integrable. Let us show that \(t \mapsto \lambda_t\) is vaguely \(\mu\) -measurable. For, we first note that if \(\pi\) and g are continuous, then the mapping \(t \mapsto \lambda_t\) is vaguely continuous. In the general case, the set of compact subsets K of T such that the restrictions of \(\pi\) and g to K are continuous is \(\mu\) -dense (Ch. IV, §5, No.~10, Prop. 15); if K is such a set, then the restriction of \(t \mapsto \lambda_t\) to K is vaguely continuous, whence the first assertion of the statement. Prop. 2 of §3, No.~1 shows that \(\Lambda\) is \(\mu\) -adequate.

Conversely, suppose that \(\Lambda\) is scalarly essentially \(\mu\) -integrable and vaguely \(\mu\) -measurable; it is then \(\mu\) -adequate (§3, No.~1, Prop. 2). Since the function 1 is lower semi-continuous on X, the function \(t \mapsto \lambda_t(1) = g(t)\) is \(\mu\) -measurable (§3, Def. 1). The set S is therefore measurable (Ch. IV, §5, No.~5, Prop. 7). The set \(\mathfrak{K}\) of compact sets K ⊂ S such that \(g|K\) is continuous and \(\Lambda|K\) is vaguely continuous is \(\mu\) -dense in S (Ch. IV, §5, No.~10, Prop. 15); if K ∈ \(\mathfrak{K}\) , then the restriction to K of the mapping \(t \mapsto \varepsilon_{\pi(t)} = \frac{1}{g(t)} \lambda_t\) is therefore vaguely continuous, and this implies the continuity of \(\pi|K\) (Ch. III, §1, No.~9, Prop. 13). Since \(\mathfrak{K}\) is \(\mu\) -dense in S, the restriction of \(\pi\) to S is \(\mu\) -measurable.

We shall make use of the following lemma:

Lemma. --- Let T and X be two topological spaces, \(\pi\) a proper continuous mapping (GT, I, §10, No.~1, Def. 1) of T into X. Let g be a lower semi-continuous numerical function defined on T. For every \(x \in X\) , let \(f(x)\) be the infimum of the function \(g(t)\) in the set \(\overline{\pi}^{-1}(x)\) (infimum equal to \(+\infty\) if \(\overline{\pi}^{-1}(x) = \varnothing\) ; cf.~S, III, §1, No.~9). Then f is lower semi-continuous on X.

For every (finite) real number \(a\) , denote by \(\mathrm{B}_a\) the set of \(x \in \mathrm{X}\) such that \(f(x) \leqslant a\) , and by \(\mathrm{A}_a\) the set of \(t \in \mathrm{T}\) such that \(g(t) \leqslant a\) ; it all comes down to showing that \(\mathrm{B}_a\) is closed (GT, IV, §6, No.~2, Prop. 1). Now, \(\mathrm{A}_a\) is closed (same ref.) and the proper mapping \(\pi\) is closed (GT, I, §10, No.~1, Prop. 1); we are thus reduced to proving that \(\pi(\mathrm{A}_a) = \mathrm{B}_a\) . The obvious relation \(f(\pi(t)) \leqslant g(t)\) for all \(t \in \mathrm{T}\) implies that \(\pi(\mathrm{A}_a) \subset \mathrm{B}_a\) . On the other hand, let \(x \in \mathrm{B}_a\) ; the set \(\overline{\pi}^{-1}(x)\) is quasi-compact (GT, I, §10, No.~2, Th. 1) and nonempty, therefore there exists a \(t \in \overline{\pi}^{-1}(x)\) such that \(g(t) = \inf_{u\in \overline{\pi}^{-1} (x)}g(u) = f(x)\) (GT, IV, §6, No.~2, Th. 3); thus \(t\in \mathbf{A}_a\) and \(\pi (t) = x\).

\subsection{Upper integrals of positive functions with respect to an integral of point measures}

We are going to see that, when \((\pi,g)\) is a \(\mu\) -adapted pair, one can sharpen the results obtained by applying the propositions of §3 to the family \(t\mapsto\lambda_{t}=g(t)\varepsilon_{\pi(t)}\) , which is \(\mu\) -adequate by Prop. 1.

THEOREM 1. --- Let \((\pi, g)\) be a \(\mu\) -adapted pair, and let

\[
\nu = \int g (t) \varepsilon_ {\pi (t)} d \mu (t).
\]

For every numerical function \(f \geqslant 0\) defined on \(X\) ,

\[
\int^ {\bullet} f (x) d \nu (x) = \int^ {\bullet} f (\pi (t)) g (t) d \mu (t).\tag{1}
\]

\begin{enumerate}
\def\labelenumi{\Alph{enumi})}
\tightlist
\item
  Suppose first that the measure \(\mu\) has compact support \(K\) and that the restrictions to \(K\) of the functions \(g\) and \(\pi\) are continuous. By formula (4) of §3, No.~1, \(\nu^{\bullet}(1) = \int_{K} g(t) d\mu(t) < +\infty\) , so that all of the measures that figure in formula (1) are bounded. We may therefore replace \(\int^{\bullet}\) by \(\int^{*}\) in the first member. In view of formula (6) of §3, No.~2, it all comes down to proving that
\end{enumerate}

\[
\int^ {*} f (x) d \nu (x) \leqslant \int^ {\bullet} f (\pi (t)) g (t) d \mu (t),\tag{2}
\]

where the symbol \(\int^{\bullet}\) in the second member can in turn be replaced by \(\int^{*}\) . By the definition of upper integral, it suffices to verify the inequality

\[
\int^ {*} f (x) d \nu (x) \leqslant \int^ {*} h (t) d \mu (t)\tag{3}
\]

for every lower semi-continuous function \(h\) on \(\mathrm{T}\) that is \(\geqslant\) the function \(t \mapsto f(\pi(t))g(t)\) . Now, let \(\varepsilon\) be a number \(> 0\) and let \(u\) be the function \((h + \varepsilon)/g\) , which is lower semi-continuous in \(\mathrm{K}\) . If \(t \in \overline{\pi}^{-1}(\{x\}) \cap \mathrm{K}\) , then \(u(t) \geqslant f(x)\) : this is obvious if \(g(t) = 0\) , because then \(u(t) = +\infty\) ; if \(g(t) > 0\) , then

\[
u (t) g (t) = h (t) + \varepsilon \geqslant f (\pi (t)) g (t) = f (x) g (t),
\]

whence the asserted inequality. Under these conditions, for every \(x \in X\) let \(v(x)\) be the infimum of \(u(t)\) for \(t \in \overline{\pi}^{-1}(\{x\}) \cap K\) . The function v is \(\geqslant f\) by the foregoing, it is lower semi-continuous on X by the Lemma (applied to the restriction of \(\pi\) to K), and \(v(\pi(t))g(t) \leqslant h(t) + \varepsilon\) for all \(t \in K\) (recall that the first member is zero by convention if \(g(t) = 0\) ). Let us then apply to v the formula (4) of §3, No.~1. We obtain:

\[
\begin{array}{l} \int^ {*} f (x) d \nu (x) \leqslant \int^ {*} v (x) d \nu (x) \\ = \int^ {*} v (\pi (t)) g (t) d \mu (t) \leqslant \int_ {K} ^ {*} (h (t) + \varepsilon) d \mu (t) \\ = \int^ {*} h (t) d \mu (t) + \varepsilon \mu (1). \end{array}\tag{4}
\]

Since the measure \(\mu\) is bounded and \(\varepsilon\) is arbitrary, the inequality (3) follows.

\begin{enumerate}
\def\labelenumi{\Alph{enumi})}
\setcounter{enumi}{1}
\tightlist
\item
  Let us now pass to the general case. Since the mapping \(t \mapsto (\pi(t), g(t))\) of \(T\) into \(X \times R_{+}\) is \(\mu\) -measurable (Ch. IV, §5, No.~3, Th. 1), the set \(\mathfrak{K}\) of compact subsets \(K\) of \(T\) such that the restrictions of \(\pi\) and \(g\) to \(K\) are continuous is \(\mu\) -dense (Ch. IV, §5, No.~10, Prop. 15). By Prop. 4 of §2, No.~3, \(\mu\) is the sum of a summable family \((\mu_{\alpha})_{\alpha \in A}\) of measures whose supports are elements of \(\mathfrak{K}\) ; the pair \((\pi, g)\) being \(\mu_{\alpha}\) -adapted for every \(\alpha \in A\) , let \(\nu_{\alpha}\) be the measure \(\int g(t)\varepsilon_{\pi(t)} d\mu_{\alpha}(t)\) . Then by A),
\end{enumerate}

\[
\int^ {\bullet} f (x) d \nu_ {\alpha} (x) = \int^ {\bullet} f (\pi (t)) g (t) d \mu_ {\alpha} (t).\tag{5}
\]

But the \(\nu_{\alpha}\) form a summable family whose sum is equal to \(\nu\) (§3, No.~1, Cor. of Prop. 1). Therefore, by Prop. 1 of §2, No.~2,

\[
\int^ {\bullet} f (x) d \nu (x) = \sum_ {\alpha \in \mathrm{A}} \int^ {\bullet} f (x) d \nu_ {\alpha} (x).\tag{6}
\]

One has an analogous relation for the second member of (5), and (1) therefore follows from (5) by summing over \(\alpha\) .

COROLLARY. --- In order that a subset N of X be locally negligible for \(\nu\) , it is necessary and sufficient that the intersection of \(\overline{\pi}^{-1}(N)\) with the set of points \(t \in T\) where \(g(t) > 0\) be locally negligible for \(\mu\) .

PROPOSITION 2. --- Let \(\pi\) be a proper continuous mapping (GT, I, §10, No.~1) of T into X, and let \(g\) be a continuous, finite numerical function on T such that \(g(t) > 0\) for all \(t \in \mathrm{T}\) . The pair \((\pi, g)\) is then \(\mu\) -adapted, and if one sets

\[
\nu = \int g (t) \varepsilon_ {\pi (t)} d \mu (t),
\]

then, for every numerical function \(f \geqslant 0\) defined on \(X\) ,

\[
\int^ {*} f (x) d \nu (x) = \int^ {*} f (\pi (t)) g (t) d \mu (t).\tag{7}
\]

It is clear that \(\pi\) and g are \(\mu\) -measurable; moreover, for every function \(\psi \in \mathcal{K}(\mathrm{X})\) , \(\psi \circ \pi\) is continuous with compact support, since \(\pi\) is proper; the pair \((\pi, g)\) is therefore \(\mu\) -adapted and, moreover, the mapping \(t \mapsto g(t) \varepsilon_{\pi(t)}\) is vaguely continuous.

Let h be a lower semi-continuous function on T such that

\[
f (\pi (t)) g (t) \leqslant h (t) \quad \text { for   all } t \in \mathrm{T}.
\]

We are going to show that

\[
\int^ {*} f (x) d \nu (x) \leqslant \int^ {*} h (t) d \mu (t).\tag{8}
\]

By the definition of upper integral, this will imply the inequality

\[
\int^ {*} f (x) d \nu (x) \leqslant \int^ {*} f (\pi (t)) g (t) d \mu (t),
\]

which, combined with the inequality (7) of §3, No.~2, will prove (7).

To prove (8), let us define a function \(\overline{f}\) on \(X\) in the following manner: \(\overline{f}(x)\) is the infimum of \(h(t) / g(t)\) in the set \(\overline{\pi}^{-1}(x)\) (infimum equal to \(+\infty\) if \(\overline{\pi}^{-1}(x) = \varnothing\) ). The function \(\overline{f}\) has the following properties:

\(1^{\circ}\overline{f} (x)\geqslant f(x)\) for all \(x\in \mathrm{X}\) (since \(g(t) > 0\) for all \(t\in \mathrm{T}\) ).

\[
2 ^ {\circ} \overline {{f}} (\pi (t)) g (t) \leqslant h (t) \text {   for   all   } t \in \mathrm{T}.
\]

\(3^{\circ}\) The function \(\overline{f}\) is lower semi-continuous by virtue of the Lemma, the function \(h / g\) being lower semi-continuous on T.

Consequently, in view of Prop. 2a) of §3, No.~1:

\[
\int^ {*} f (x) d \nu (x) \leqslant \int^ {*} \overline {{{f}}} (x) d \nu (x) = \int^ {*} \overline {{{f}}} (\pi (t)) g (t) d \mu (t) \leqslant \int^ {*} h (t) d \mu (t),
\]

which establishes (8), and completes the proof.

\subsection{Measurability with respect to an integral of point measures}

PROPOSITION 3. --- Let \((\pi, g)\) be a \(\mu\) -adapted pair, and let

\[
\nu = \int g (t) \varepsilon_ {\pi (t)} d \mu (t).
\]

Let \(f\) be a mapping of \(X\) into a topological space \(G\) , and let \(S\) be the \((\mu\text{-measurable})\) set of points \(t \in T\) such that \(g(t) > 0\) . In order that \(f\) be \(\nu\) -measurable, it is necessary and sufficient that the restriction of \(f \circ \pi\) to \(S\) be \(\mu\) -measurable.

Suppose first that \(f\) is \(\nu\) -measurable. By hypothesis, the set \(\mathcal{K}\) of compact subsets \(K\) of \(S\) , such that the restriction of \(\pi\) to \(K\) is continuous, is \(\mu\) -dense in \(S\) (Ch. IV, §5, No.~10, Prop. 15). To show that the restriction of \(f \circ \pi\) to \(S\) is \(\mu\) -measurable, it therefore suffices to prove that for every \(K \in \mathcal{K}\) , the set of compact subsets \(H\) of \(K\) , such that the restriction of \(f \circ \pi\) to \(H\) is continuous, is \(\mu\) -dense in \(K\) (Ch. IV, §5, No.~8, Prop. 13). But by hypothesis, there exists a partition of the compact set \(\pi(K)\) formed by a \(\nu\) -negligible set \(N\) and a sequence of compact sets \((C_n)\) such that the restriction of \(f\) to each \(C_n\) is continuous. Under these conditions, \(K \cap \overline{\pi}^{-1}(N)\) and the sets \(K \cap \overline{\pi}^{-1}(C_n)\) form a partition of \(K\) ; but \(K \cap \overline{\pi}^{-1}(N)\) is \(\mu\) -negligible by virtue of the Cor. of Th. 1 of No.~2, the sets \(K \cap \overline{\pi}^{-1}(C_n)\) are compact, and the restriction of \(f \circ \pi\) to each of the latter sets is continuous, which proves that the restriction of \(f \circ \pi\) to \(S\) is \(\mu\) -measurable.

Conversely, suppose this to be the case; to show that \(f\) is \(\nu\) -measurable, it suffices to prove that the set \(\mathfrak{L}\) of compact subsets \(L\) of \(X\) , such that the restriction of \(f\) to \(L\) is continuous, is \(\nu\) -dense (Ch. IV, §5, No.~10, Prop. 15). Let \(N\) be a subset of \(X\) such that \(N \cap L\) is \(\nu\) -negligible for every \(L \in \mathfrak{L}\) , and let us show that \(N\) is locally \(\nu\) -negligible. For this, we must show that \(\overline{\pi}^{-1}(N) \cap S\) is locally \(\mu\) -negligible (Cor. of Th. 1 of No.~2). Now, the set \(\mathfrak{H}\) of compact subsets \(H\) of \(S\) , such that the restrictions to \(H\) of \(\pi\) and \(f \circ \pi\) are continuous, is by hypothesis \(\mu\) -dense in \(S\) (Ch. IV, §5, No.~10, Prop. 15). It therefore suffices to prove that \(\overline{\pi}^{-1}(N) \cap H\) is \(\mu\) -negligible for every \(H \in \mathfrak{H}\) . Now, \(\pi(H)\) is compact and may be identified with the quotient space of \(H\) by the equivalence relation \(\pi(t) = \pi(t')\) , \(\pi\) being identified with the canonical mapping of \(H\) onto this quotient space (GT, I, §5, No.~2, Prop. 3). Since the restriction of \(f \circ \pi\) to \(H\) is continuous, the restriction of \(f\) to \(\pi(H)\) is therefore continuous, in other words \(\pi(H) \in \mathfrak{L}\) , consequently \(N \cap \pi(H)\) is \(\nu\) -negligible. By the Cor. of Th. 1 of No.~2, \(\overline{\pi}^{-1}(N \cap \pi(H)) \cap S\) is locally \(\mu\) -negligible; the same is therefore true of the set

\[
\mathrm{H} \cap \overline {{\pi}} ^ {- 1} (\mathrm{N}) \subset \overline {{\pi}} ^ {- 1} \bigl (\mathrm{N} \cap \pi (\mathrm{H}) \bigr) \cap \mathrm{S};
\]

but since H is compact, \(\mathrm{H} \cap \overline{\pi}^{-1}(\mathrm{N})\) is \(\mu\) -negligible, which completes the proof.

Remark. --- If f is a mapping of X into a Banach space F, it comes to the same to say that the restriction of \(f \circ \pi\) to S is \(\mu\) -measurable or to say that the function \((\mathbf{f} \circ \pi)g\) (defined on T) is \(\mu\) -measurable, since \(g\) is \(\mu\) -measurable, is not zero in S, and is zero on T - S (Ch. IV, §5, No.~10, Prop. 15).

\subsection{Integration of functions with values in a Banach space, with respect to an integral of point measures}

THEOREM 2. --- Let \((\pi, g)\) be a \(\mu\) -adapted pair, and let

\[
\nu = \int g (t) \varepsilon_ {\pi (t)} d \mu (t).
\]

Let \(\mathbf{f}\) be a function defined on \(\mathbf{X}\) , with values in a Banach space \(\mathbf{F}\) or in \(\overline{\mathbf{R}}\) . For \(\mathbf{f}\) to be essentially \(\nu\) -integrable, it is necessary and sufficient that \(t \mapsto \mathbf{f}(\pi(t))g(t)\) be essentially \(\mu\) -integrable, in which case

\[
\int \mathbf {f} (x) d \nu (x) = \int \mathbf {f} (\pi (t)) g (t) d \mu (t).\tag{9}
\]

Suppose, moreover, that \(\pi\) is continuous and proper, and that g is continuous and such that \(g(t) > 0\) for all \(t \in T\) . Then, for f to be \(\nu\) -integrable, it is necessary and sufficient that \(t \mapsto \mathbf{f}(\pi(t))g(t)\) be \(\mu\) -integrable.

\begin{enumerate}
\def\labelenumi{\Alph{enumi})}
\tightlist
\item
  We begin by treating the case that the measure \(\mu\) has compact support K, on which g is bounded. The measures \(\mu\) and \(\nu\) are then bounded, and one can replace `essentially integrable' in the statement by `integrable'. Suppose that f is \(\nu\) -integrable: the function \(\mathbf{f}(\pi(t))g(t)\) is then \(\mu\) -integrable, and the relation (9) is verified, by Th. 1 of §3, No.~3. Conversely, suppose that \(\mathbf{f}(\pi(t))g(t)\) is \(\mu\) -integrable: f is then \(\nu\) -measurable (No.~3, Prop. 3 and Remark), and
\end{enumerate}

\[
\int^ {\bullet} | \mathbf {f} (x) | d \nu (x) = \int^ {\bullet} | \mathbf {f} (\pi (t)) | g (t) d \mu (t) <   + \infty
\]

(No.~2, Th. 1); \(\mathbf{f}\) is therefore essentially \(\nu\) -integrable (§1, No.~3, Prop. 9), hence \(\nu\) -integrable. Th. 1 of §3, No.~3 then implies (9).

\begin{enumerate}
\def\labelenumi{\Alph{enumi})}
\setcounter{enumi}{1}
\tightlist
\item
  Let us pass to the general case. Let \(\mathfrak{K}\) be the set of compact subsets \(K\) of \(T\) such that \(g|K\) is continuous: \(\mathfrak{K}\) is \(\mu\) -dense (Ch. IV, §5, No.~10, Prop. 15), therefore the measure \(\mu\) is the sum of a family \((\mu_{\alpha})_{\alpha \in A}\) of measures whose supports are elements of \(\mathfrak{K}\) (§2, No.~3, Prop. 4). The pair \((g,\pi)\) is obviously \(\mu_{\alpha}\) -adapted for every \(\alpha \in A\) , and the measure \(\nu\) is the sum of the family of measures \(\nu_{\alpha} = \int \varepsilon_{\pi(t)} g(t) d\mu_{\alpha}(t)\) (§3, No.~1, Cor. of Prop. 1). Since the argument of A) may be applied to the measures \(\mu_{\alpha}, \nu_{\alpha}\) , the first part of the statement then follows from Prop. 3 of §2, No.~2.
\end{enumerate}

For the function \(\mathbf{f}\) (resp. \(t\mapsto \mathbf{f}(\pi (t))g(t)\) ) to be integrable for \(\nu\) (resp. for \(\mu\) ), it is necessary and sufficient that it be essentially integrable and that

\[
\int^ {*} | \mathbf {f} (x) | d \nu (x) <   + \infty \quad (\mathrm{resp.} \int^ {*} | \mathbf {f} (\pi (t)) | g (t) d \mu (t) <   + \infty).
\]

The second part of the statement therefore follows from the first part and Proposition 2.

Remark. --- Let \((\pi, g)\) be a \(\mu\) -adapted pair, \(\pi'\) a mapping of T into X, and \(g'\) a finite numerical function \(\geqslant 0\) defined on T, such that \(\pi'\) (resp. \(g'\) ) is equal to \(\pi\) (resp. g) locally almost everywhere for \(\mu\) . Then the pair \((\pi', g')\) is \(\mu\) -adapted, the measures \(\lambda_t = g(t)\varepsilon_{\pi(t)}\) and \(\lambda_t' = g'(t)\varepsilon_{\pi'(t)}\) are equal locally almost everywhere, and \(\int g(t)\varepsilon_{\pi(t)} d\mu(t) = \int g'(t)\varepsilon_{\pi'(t)} d\mu(t)\) . If now \(\pi'\) and \(g'\) are only defined locally almost everywhere (for \(\mu\) ) and if there exists a \(\mu\) -adapted pair \((\pi, g)\) such that \(\pi'\) (resp. \(g'\) ) is equal to \(\pi\) (resp. g) locally almost everywhere, one again says that the pair \((\pi', g')\) is \(\mu\) -adapted and one sets

\[
\int g ^ {\prime} (t) \varepsilon_ {\pi^ {\prime} (t)} d \mu (t) = \int g (t) \varepsilon_ {\pi (t)} d \mu (t)
\]

(cf.~§3, No.~3, Remark). The statements of Ths. 1 and 2 and of Prop. 3 remain valid when \(\pi\) and \(g\) are only assumed to be defined locally almost everywhere.

\section{MEASURES DEFINED BY NUMERICAL DENSITIES}

\subsection{Locally integrable functions}

PROPOSITION 1. --- Let g be a function defined locally almost everywhere in T (for the positive measure \(\mu\) ), with values in a Banach space F (resp. in \(\overline{R}\) ). The following properties are equivalent:

\begin{enumerate}
\def\labelenumi{\alph{enumi})}
\item
  For every point \(t \in \mathbb{T}\) , there exists a neighborhood \(\mathbb{V}\) of \(t\) such that the function \(\mathbf{g}\varphi_{\mathbb{V}}\) is \(\mu\) -integrable.
\item
  The function \(\mathbf{g}\) is \(\mu\) -measurable and, for every compact set \(\mathrm{K} \subset \mathrm{T}\) , \(\int^{*} |\mathbf{g}| \varphi_{\mathrm{K}} d\mu < +\infty\) .
\item
  For every numerical function \(h \in \mathcal{K}(\mathrm{T})\) , gh is \(\mu\) -integrable.
\end{enumerate}

Let us show that a) implies b); for, the function g is measurable by the principle of localization (Ch. IV, §5, No.~2, Prop. 4). On the other hand, for every \(t \in \mathbf{K}\) there exists, by hypothesis, a neighborhood \(V_t\) of \(t\) in \(T\) such that \(\mathbf{g}\varphi_{V_t}\) is integrable; one can therefore cover \(K\) by a finite number of neighborhoods \(V_i\) ( \(1 \leqslant i \leqslant n\) ) such that the functions \(\mathbf{g}\varphi_{V_i}\) are integrable. Since \(|\mathbf{g}| \varphi_K \leqslant \sum_{i=1}^{n} |\mathbf{g}| \varphi_{V_i}\) , one has \(\int^* |\mathbf{g}| \varphi_K d\mu < +\infty\) .

Secondly, b) implies c), because gh is then measurable, and if L is the compact support of h then \(|gh| \leqslant \|h\|\cdot|g|\varphi_{L}\) , therefore \(\int^{*}|gh|d\mu < +\infty\) by hypothesis; gh is therefore integrable by the criterion for integrability (Ch. IV, §5, No.~6, Th. 5).

Finally, c) implies a). Indeed, for every \(t \in T\) let V be a compact neighborhood of t. There exists a continuous mapping h of T into {[}0,1{]}, equal to 1 on V and with compact support (Ch. III, §1, No.~2, Lemma 1); by hypothesis gh is integrable, therefore so is \(\mathbf{g}\varphi_{\mathrm{V}} = (\mathbf{g}h)\varphi_{\mathrm{V}}\) (Ch. IV, §5, No.~6, Cor. 3 of Th. 5).

DEFINITION 1. --- A function g, defined locally almost everywhere in T (for the positive measure \(\mu\) ), with values in a Banach space F (resp. in \(\overline{R}\) ), is said to be locally integrable for \(\mu\) (or locally \(\mu\) -integrable) if it satisfies the conditions a), b), c) of Prop. 1. If \(\theta\) is a complex measure, a function g defined locally \(\theta\) -almost everywhere is said to be locally \(\theta\) -integrable if it is locally integrable for the positive measure \(|\theta|\) .

If \(\mathbf{g}\) is locally \(\theta\) -integrable, then every function equal to \(\mathbf{g}\) locally almost everywhere is locally integrable. It is clear that the sum of two locally integrable functions is locally integrable. The functions with values in \(\mathbf{F}\) , everywhere defined and locally integrable for \(\theta\) , form a vector space denoted \(\mathcal{L}_{\mathrm{loc}}^{1}(\mathrm{T},\theta ;\mathrm{F})\) ; when \(\mathbf{F} = \mathbf{R}\) or \(\mathbf{C}\) , the mention of \(\mathbf{F}\) is often omitted if there is no ambiguity. This space will always be equipped (absent express mention to the contrary) with the topology defined by the semi-norms \(\mathbf{g}\mapsto \int |\mathbf{g}\varphi_{\mathbf{K}}|d|\theta |\) , where \(\mathbf{K}\) runs over the set of compact subsets of \(\mathbf{T}\) . The associated Hausdorff space, the quotient of \(\mathcal{L}_{\mathrm{loc}}^{1}(\mathrm{T},\theta ;\mathrm{F})\) by the subspace \(\mathcal{N}_{\mathrm{F}}^{\infty}\) of mappings that are zero locally almost everywhere, is denoted \(\mathrm{L}_{\mathrm{loc}}^{1}(\mathrm{T},\theta ;\mathrm{F})\) . The spaces \(\mathrm{L}_{\mathrm{loc}}^{1}(\mathrm{T},\theta ;\mathrm{F})\) and \(\mathrm{L}_{\mathrm{loc}}^{1}(\mathrm{T},|\theta |; \mathrm{F})\) are identical.

It can be shown that the topological vector spaces just defined are complete (Exer. 31).

Every measurable function g, that is essentially bounded on every compact set, is locally integrable. For every number p such that \(1 \leqslant p \leqslant +\infty\) , every function \(g \in L_{F}^{p}\) is locally integrable; indeed, for every function \(h \in \mathcal{K}(T)\) , h belongs to \(L^{q}\) (where q is the exponent conjugate to p), therefore gh is integrable (Ch. IV, §6, No.~4, Cor. 4 of Th. 2).

Let \(F, G, H\) be three Banach spaces, and \((\mathbf{u}, \mathbf{v}) \mapsto \Phi(\mathbf{u}, \mathbf{v})\) a continuous bilinear mapping of \(F \times G\) into \(H\) . If \(f\) is locally integrable and takes its values in F, and if \(g \in L_{G}^{\infty}\) , then \(\Phi(\mathbf{f}, \mathbf{g})\) is locally integrable (Ch. IV, §6, No.~4, Cor. 1 of Th. 2).

\subsection{Measures defined by numerical densities}

Let \(g\) be a positive numerical function defined locally \(\mu\) -almost everywhere in \(T\) and locally \(\mu\) -integrable; the set of \(t\) such that \(g(t) = +\infty\) is then locally \(\mu\) -negligible, because \(g\varphi_{K}\) is \(\mu\) -integrable for every compact set \(K\) (Ch. IV, §2, No.~3, Prop. 7). Now let \(g'\) be a locally integrable function that is positive and finite, equal to \(g\) locally \(\mu\) -almost everywhere; set \(\lambda_t' = g'(t)\varepsilon_t\) . The mapping \(t \mapsto \lambda_t'\) of \(T\) into \(\mathcal{M}_+(T)\) is vaguely \(\mu\) -measurable and scalarly essentially integrable (or again, the pair \((I,g')\) , where I is the identity mapping of \(T\) , is \(\mu\) -adapted); the integral \(\nu = \int \lambda_t' d\mu(t)\) does not depend on the particular function \(g'\) , locally almost everywhere equal to \(g\) , used in the definition of the measures \(\lambda_t'\) . This measure \(\nu\) is determined by the condition

\[
\int f (t) d \nu (t) = \int f (t) g (t) d \mu (t) \quad \mathrm{for} f \in \mathcal {K} (\mathrm{T}).\tag{1}
\]

If now \(\theta\) is a complex measure, and if u is a complex function (or a function with values in \(\overline{R}\) ) defined locally \(\theta\) -almost everywhere and locally integrable for \(\theta\) , one can write

\[
\begin{array}{l} {u = g _ {1} - g _ {2} + i (g _ {3} - g _ {4})} \\ {\theta = \mu_ {1} - \mu_ {2} + i (\mu_ {3} - \mu_ {4})} \end{array}\tag{2}
\]

where \(\mu_{1} = (\mathcal{R}\theta)^{+}\) , \(\mu_{2} = (\mathcal{R}\theta)^{-}\) , \(\mu_{3} = (\mathcal{I}\theta)^{+}\) , \(\mu_{4} = (\mathcal{I}\theta)^{-}\) (Ch. III, §1, No.~5), and where \(g_{1}, g_{2}, g_{3}, g_{4}\) have the analogous meanings; since \(|u|\) is locally \(|\theta|\) -integrable, each of the positive functions \(g_{i}\) ( \(i = 1, 2, 3, 4\) ) is locally integrable for each positive measure \(\mu_{j}\) ( \(j = 1, 2, 3, 4\) ), so that the mapping

\[
f \mapsto \int f (t) u (t) d \theta (t)
\]

on \(\mathcal{K}(\mathrm{T})\) is a complex measure.

DEFINITION 2. --- Let \(\theta\) be a complex measure, and let \(u\) be a complex function (or a function with values in \(\overline{\mathbf{R}}\) ) defined locally \(\theta\) -almost everywhere and locally \(\theta\) -integrable. The complex measure \(f \mapsto \int f u d\theta\) on T is said to be the product of the measure \(\theta\) by the function \(u\) , or the measure with density \(u\) with respect to \(\theta\) , and is denoted \(u \cdot \theta\) .

Every complex measure that is the product of a positive measure \(\mu\) by a locally \(\mu\) -integrable function is called a measure with base \(\mu\) .

The relation \(\eta = u \cdot \theta\) is again, by convention, written

\[
d \eta (t) = u (t) d \theta (t).
\]

When u is everywhere defined and continuous, one recovers the definition given in Ch. III, §1, No.~4. It is clear that if \(u_{1}\) and \(u_{2}\) are locally \(\theta\) -integrable, then \((u_{1}+u_{2})\cdot\theta=u_{1}\cdot\theta+u_{2}\cdot\theta\) . Similarly, if \(\theta_{1}\) and \(\theta_{2}\) are two measures on T, and if u is a function locally integrable for \(\theta_{1}\) and \(\theta_{2}\) , then u is locally integrable for \(\theta_{1}+\theta_{2}\) and one has \(u\cdot(\theta_{1}+\theta_{2})=u\cdot\theta_{1}+u\cdot\theta_{2}\) .

We shall henceforth restrict ourselves to the case of functions defined everywhere; the extension to functions defined locally almost everywhere, which is always obvious, is left to the reader.

The following proposition permits, for the most part, reducing the case of complex measures to that of positive measures:

PROPOSITION 2. --- Let \(\theta\) be a complex measure, and \(u\) a locally \(\theta\) -integrable complex function; then

\[
\left| u \cdot \theta \right| = \left| u \right| \cdot \left| \theta \right|.\tag{3}
\]

We begin with an auxiliary result:

Lemma 1. --- Let \(\theta\) be a complex measure, and let \(f\) be an element of \(\overline{\mathcal{L}}_{\mathbf{C}}^{1}(\mathrm{T},\theta)\) . Then

\[
\langle | \theta |, | f | \rangle = \sup _ {c \in \mathcal {K} _ {1}} | \langle \theta , c f \rangle | = \sup _ {c \in \mathcal {B} _ {1}} | \langle \theta , c f \rangle |,\tag{4}
\]

where \(\mathcal{K}_{1}\) (resp. \(B_{1}\) ) denotes the set of complex functions c, continuous with compact support (resp. Borel), such that \(|c| \leqslant 1\) .

Let us first treat the case that \(f \in \mathcal{K}(\mathrm{T};\mathbf{C})\) . Obviously

\[
\sup _ {c \in \mathcal {K} _ {1}} | \langle \theta , c f \rangle | \leqslant \sup _ {c \in \mathcal {B} _ {1}} | \langle \theta , c f \rangle | \leqslant \langle | \theta |, | f | \rangle
\]

(Ch. IV, §4, No.~2, Prop. 2). On the other hand, let \(g\) be an element of \(\mathcal{K}(\mathrm{T};\mathbf{C})\) such that \(|g| \leqslant |f|\) ; \(g\) is the uniform limit of a sequence \((g_n)\) of elements of \(\mathcal{K}(\mathrm{T};\mathbf{C})\) whose supports are contained in the open set \(\mathrm{U}\) formed by the \(t\) such that \(f(t) \neq 0\) , and one may clearly suppose that \(|g_n| \leqslant |f|\) for every \(n\) . Set \(c_n(t) = g_n(t)/f(t)\) for \(t \in \mathrm{U}\) , \(c_n(t) = 0\) for \(t \notin \mathrm{U}\) ; then \(c_n \in \mathcal{K}_1\) , \(g = \lim_{n \to \infty} c_n f\) , therefore \(|\langle \theta, g \rangle| = \lim_{n \to \infty} |\langle \theta, c_n f \rangle|\) , and finally

\[
\sup _ {| g | \leqslant | f |, g \in \mathscr {K} (\mathrm{T}; \mathbf {C})} | \langle \theta , g \rangle | \leqslant \sup _ {c \in \mathscr {K} _ {1}} | \langle \theta , c f \rangle |.
\]

One concludes by observing that the first member of this inequality is equal to \(\langle|\theta|,|f|\rangle\) (Ch. III, §1, No.~6, formula (12)).

Next, denote by \(f\) an element of \(\overline{\mathcal{L}}_{\mathbf{C}}^{1}(\theta)\) , and let us show that (4) is again true: it suffices to verify that the three members of this relation depend continuously on \(f\) for the topology of \(\overline{\mathcal{L}}_{\mathbf{C}}^{1}(\theta)\) , since they coincide on the dense subspace \(\mathcal{K}(\mathrm{T};\mathbf{C})\) . This results at once from the following inequalities, where \(f\) and \(f'\) denote elements of \(\overline{\mathcal{L}}_{\mathbf{C}}^{1}(\theta)\) :

\[
| \langle | \theta |, | f | \rangle - \langle | \theta |, | f ^ {\prime} | \rangle | \leqslant \langle | \theta |, | f - f ^ {\prime} | \rangle = \overline {{\mathrm{N}}} _ {1} (f - f ^ {\prime})
\]

\[
| \langle \theta , c f \rangle - \langle \theta , c f ^ {\prime} \rangle | \leqslant \langle | \theta |, | c | | f - f ^ {\prime} | \rangle \leqslant \overline {{{\mathrm{N}}}} _ {1} (f - f ^ {\prime})
\]

for all \(c \in B_{1}\) . The lemma is thus established.

Passing to the proof of Proposition 2, let us apply the lemma to the function uf, where h belongs to \(\mathcal{K}_{+}(\mathrm{T})\) . This yields:

\[
\langle | \theta |, | u h | \rangle = \sup _ {c \in \mathscr {K} _ {1}} | \langle \theta , c u h \rangle | = \sup _ {c \in \mathscr {K} _ {1}} | \langle u \cdot \theta , c h \rangle | = \langle | u \cdot \theta |, h \rangle .
\]

However, the first member is also equal to

\[
\langle | \theta |, | u | h \rangle = \langle | u | \cdot | \theta |, h \rangle .
\]

The two measures \(|u| \cdot |\theta|\) and \(|u \cdot \theta|\) are therefore equal.

COROLLARY. --- Let \(g_{1}\) and \(g_{2}\) be two locally \(\mu\) -integrable numerical functions; then

\[
\inf (g _ {1} \cdot \mu , g _ {2} \cdot \mu) = \inf (g _ {1}, g _ {2}) \cdot \mu ; \quad \sup (g _ {1} \cdot \mu , g _ {2} \cdot \mu) = \sup (g _ {1}, g _ {2}) \cdot \mu .
\]

In particular, if \(g\) is a locally \(\mu\) -integrable numerical function, then

\[
(g \cdot \mu) ^ {+} = g ^ {+} \cdot \mu ; (g \cdot \mu) ^ {-} = g ^ {-} \cdot \mu .
\]

This follows at once from Prop. 2 and the formulas (6) of Ch. II, §1, No.~1.

\subsection{Integration with respect to a measure defined by a density}

In the statements of this subsection, g denotes a positive numerical function, defined everywhere and locally \(\mu\) -integrable, \(\theta\) denotes a complex measure, and u a locally \(\theta\) -integrable complex function.

The remarks in No.~2 show that the results of §4 are applicable to the measure \(\nu = g \cdot \mu = \int g(t)\varepsilon_t d\mu(t)\) (even though the measure \(g(t)\varepsilon_t\) is not defined unless \(g(t) \neq +\infty\) ). We thus obtain the following statement:

PROPOSITION 3. --- For every numerical function \(f \geqslant 0\) defined on T,

\[
\int^ {\bullet} f d \nu = \int^ {\bullet} (f g) d \mu .\tag{5}
\]

This follows from Th. 1 of §4, No.~2.

COROLLARY 1. --- In order that a function f, with values in a Banach space or in \(\overline{R}\) , be locally negligible for the measure \(u \cdot \theta\) , it is necessary and sufficient that uf be locally negligible for \(\theta\) .

To say that f (resp. \(uf\) ) is locally negligible for \(u \cdot \theta\) (resp. for \(\theta\) ) is equivalent to saying that \(|f|\) (resp. \(|u| \cdot |f|\) ) is locally negligible for \(|u \cdot \theta|\) (resp. for \(|\theta|\) ). We are thus reduced, in view of Prop. 2 of No.~2, to the case that \(f, u, \theta\) are positive; the statement then follows at once from Prop. 3.

COROLLARY 2. --- Let \(u_1\) and \(u_2\) be two locally \(\theta\) -integrable complex functions. In order that \(u_1 \cdot \theta = u_2 \cdot \theta\) , it is necessary and sufficient that \(u_1\) and \(u_2\) be equal locally almost everywhere.

One is immediately reduced to showing that \(u \cdot \theta = 0\) implies that \(u(t) = 0\) locally almost everywhere; but \(u \cdot \theta = 0\) means that the function 1 is locally negligible for the measure \(u \cdot \theta\) . One then applies Corollary 1.

COROLLARY 3. --- Let u be a complex function that is locally integrable for the positive measure \(\mu\) . For \(u \cdot \mu\) to be a positive measure, it is necessary and sufficient that \(u(t) \geqslant 0\) locally almost everywhere.

For, \(u \cdot \mu\) is positive if and only if \(u \cdot \mu = |u \cdot \mu| = |u| \cdot \mu\) (Prop. 2), and this is equivalent to \(u = |u|\) locally almost everywhere (Cor. 2).

PROPOSITION 4. --- For a mapping f of T into a topological space G to be measurable for the measure \(u \cdot \theta\) , it is necessary and sufficient that the restriction of f, to the \(\theta\) -measurable set S of the t such that \(u(t) \neq 0\) , be \(\theta\) -measurable.

When u and \(\theta\) are positive, this follows at once from Prop. 3 of §4, No.~3. The result then extends to the case that u and \(\theta\) are complex thanks to Prop. 2.

COROLLARY. --- Let f be a function defined on T, with values in a Banach space F or in \(\overline{R}\) . For f to be \((u\cdot\theta)\) -measurable, it is necessary and sufficient that uf be \(\theta\) -measurable.

For, \(u\mathbf{f}\) is the extension by 0 of \((uf)|S\) to \(T\) .

THEOREM 1. --- Let f be a function defined on T, with values in a Banach space F or in \(\overline{R}\) . In order that f be essentially integrable for the measure \(\eta = u\cdot \theta\) , it is necessary and sufficient that uf be essentially \(\theta\) -integrable, in which case

\[
\int \mathbf {f} d \eta = \int (u \mathbf {f}) d \theta .\tag{6}
\]

Suppose moreover that \(u\) is continuous and that \(u(t) \neq 0\) for all \(t \in \mathbb{T}\) ; then \(\mathbf{f}\) is integrable for the measure \(\eta\) if and only if \(u\mathbf{f}\) is integrable for \(\theta\) .

The case that u and \(\theta\) are positive follows at once from Th. 2 of §4, No.~4. The first and the last assertion of the statement then follow immediately, because f is essentially integrable (resp. integrable) with respect to \(\eta = u \cdot \theta\) if and only if it is essentially integrable (resp. integrable) for \(|\eta| = |u| \cdot |\theta|\) . Finally, suppose that f is essentially integrable for \(\eta\) (hence for \(|\eta|\) ); we make use of the decomposition (2): f is essentially integrable for each of the measures \(\eta_{ij} = g_i \cdot \mu_j\) (i = 1, 2, 3, 4, j = 1, 2, 3, 4), because these measures are \(\leqslant |\eta|\) . We have

\[
\int \mathbf {f} d \eta_ {i j} = \int g _ {i} \mathbf {f} d \mu_ {j}.
\]

The formula (6) follows immediately from this.

COROLLARY. --- For the measure \(u \cdot \theta\) to be bounded, it is necessary and sufficient that u be essentially \(\theta\) -integrable.

Example. --- Let A be a subset of T; for \(\varphi_{A}\) to be locally \(\mu\) -integrable, it is necessary and sufficient that A be \(\mu\) -measurable. Assuming the condition to be fulfilled, set \(\nu = \varphi_{A} \cdot \mu\) ; for every numerical function \(f \geqslant 0\) defined on T, we then have

\[
\int^ {\bullet} f d \nu = \int^ {\bullet} f \varphi_ {\mathrm{A}} d \mu ,
\]

a quantity that is also denoted by \(\int_{\mathrm{A}}^{\bullet}f d\mu\) (cf.~Ch. IV, §5, No.~6). For a mapping \(g\) of T into a topological space G to be \(\nu\) -measurable, it is necessary and sufficient that the restriction of \(g\) to A be \(\mu\) -measurable. For a mapping \(\mathbf{f}\) of T into a Banach space F or into \(\overline{\mathbf{R}}\) to be essentially \(\nu\) -integrable, it is necessary and sufficient that \(\mathbf{f}\varphi_{\mathrm{A}}\) be essentially \(\mu\) -integrable, in which case

\[
\int \mathbf {f} d \nu = \int \mathbf {f} \varphi_ {\mathrm{A}} d \mu ,
\]

an expression that is also denoted \(\int_{A} f d\mu\) . Note that if two mappings of T into G (resp. F, \(\overline{R}\) ) coincide on A, then for one of them to be \(\nu\) -measurable (resp. essentially \(\nu\) -integrable), it is necessary and sufficient that the other be so. If now g is a mapping into G of an arbitrary subset \(B \supset A\) of T, one says that g is \(\mu\) -measurable on A if an arbitrary extension to T of the restriction of g to A is \(\nu\) -measurable, which amounts to saying that the restriction of g to A is \(\mu\) -measurable. A mapping f of B into a Banach space F, or into \(\overline{R}\) , is said to be essentially \(\mu\) -integrable on A if some extension \(\overline{f}\) to T of the restriction of f to A is essentially \(\nu\) -integrable; one then sets

\[
\int_ {\mathrm{A}} \mathbf {f} d \mu = \int_ {\mathrm{A}} \overline {{\mathbf {f}}} d \mu = \int \overline {{\mathbf {f}}} \varphi_ {\mathrm{A}} d \mu ,
\]

and one says that \(\int_{\mathrm{A}} \mathbf{f} d\mu\) is the integral of \(\mathbf{f}\) on \(\mathrm{A}\) (or extended to \(\mathrm{A}\) ). If \(f\) is a numerical function \(\geqslant 0\) defined on \(\mathrm{B} \supset \mathrm{A}\) , one defines similarly \(\int_{\mathrm{A}}^{*} f d\mu\) and \(\int_{\mathrm{A}}^{\bullet} f d\mu\) . Finally, a numerical function \(g\) defined on \(\mathrm{B} \supset \mathrm{A}\) is said to be locally \(\mu\) -integrable on \(\mathrm{A}\) if some extension \(\overline{g}\) to \(\mathrm{T}\) of the restriction of \(g\) to \(\mathrm{A}\) is locally \(\nu\) -integrable: this is equivalent to saying that, for every compact subset \(\mathrm{K}\) of \(\mathrm{T}\) , \(\overline{g}\varphi_{\mathrm{K} \cap \mathrm{A}}\) is \(\mu\) -integrable.

\subsection{Behavior of the product with respect to the usual operations}

PROPOSITION 5. --- Let \((\lambda_{\alpha})_{\alpha\in\mathbb{A}}\) be a family of positive measures on T, directed for the relation \(\leqslant\) , admitting in \(\mathcal{M}(\mathrm{T})\) a supremum \(\lambda\) . For a positive numerical function g to be locally \(\lambda\) -integrable, it is necessary and sufficient that g be locally \(\lambda_{\alpha}\) -integrable for every \(\alpha\in A\) and that the family \((g\cdot\lambda_{\alpha})\) be bounded above in \(\mathcal{M}(\mathrm{T})\) ; in this case,

\[
g \cdot \lambda = \sup _ {\alpha \in \mathrm{A}} g \cdot \lambda_ {\alpha}.
\]

It is clear that the condition is necessary. Conversely, suppose that \(g\) is locally integrable for each of the measures \(\lambda_{\alpha}\) and that the family \((g \cdot \lambda_{\alpha})_{\alpha \in \mathbb{A}}\) is bounded above; denote its supremum by \(\lambda'\) . The function \(g\) is then \(\lambda\) -measurable (§1, No.~4, Cor. 2 of Prop. 11); moreover, for every function \(h \in \mathcal{K}_{+}(\mathbf{T})\) ,

\[
\int^ {\bullet} (h g) d \lambda = \sup _ {\alpha \in A} \int^ {\bullet} (h g) d \lambda_ {\alpha} = \sup _ {\alpha \in A} \int^ {\bullet} h d (g \cdot \lambda_ {\alpha}) = \int^ {\bullet} h d \lambda^ {\prime}
\]

(§1, No.~4, Prop. 11). This implies first of all that the first member is finite for any \(h\) , so that \(g\) is locally \(\lambda\) -integrable; the symbol \(\int^{\bullet}\) may therefore be replaced by \(\int\) , and the formula may be written \(\int h d(g \cdot \lambda) = \int h d\lambda'\) . It follows that \(g \cdot \lambda = \lambda'\) , and this completes the proof.

COROLLARY. --- Suppose that \(\mu\) is the sum of a family \((\mu_{\alpha})_{\alpha\in\mathrm{A}}\) of measures on T. For a positive numerical function g defined on T to be locally \(\mu\) -integrable, it is necessary and sufficient that g be locally \(\mu_{\alpha}\) -integrable for every \(\alpha\in A\) and that the family \((g\cdot\mu_{\alpha})_{\alpha\in\mathrm{A}}\) be summable. In this case,

\[
g \cdot \mu = \sum_ {\alpha \in A} g \cdot \mu_ {\alpha}.\tag{7}
\]

Let \((g_{\alpha})_{\alpha \in A}\) be a family of \(\mu\) -measurable positive functions defined on T. Let \(S_{\alpha}\) be the set of \(t \in T\) such that \(g_{\alpha}(t) \neq 0\) . We shall say that the family \((g_{\alpha})\) is locally countable if the family \((S_{\alpha})\) is locally countable (Ch. IV, §5, No.~9); this amounts to saying that, for every compact set K in T, the set of \(\alpha \in A\) such that \(g_{\alpha}|K\) is not zero is countable.

PROPOSITION 6. --- Let \((g_{\alpha})_{\alpha \in A}\) be a locally countable family of locally \(\mu\) -integrable positive numerical functions defined on T. In order that the function \(g = \sum_{\alpha \in A} g_{\alpha}\) be locally \(\mu\) -integrable, it is necessary and sufficient that the family of measures \((g_{\alpha} \cdot \mu)_{\alpha \in A}\) be summable, in which case

\[
g \cdot \mu = \sum_ {\alpha \in A} g _ {\alpha} \cdot \mu .\tag{8}
\]

It is clear that \(g\) is \(\mu\) -measurable (Ch. IV, §5, No.~2, Prop. 4 and No.~4, Cor. 1 of Th. 2). For \(g\) to be locally \(\mu\) -integrable, it is therefore necessary and sufficient that \(\mu^{\bullet}(gf)\) be finite for every \(f \in \mathcal{K}_{+}(\mathrm{T})\) . Now, since the set of \(\alpha \in \mathrm{A}\) such that \(g_{\alpha}f \neq 0\) is countable, we have \(\mu^{\bullet}(gf) = \sum_{\alpha \in \mathrm{A}} \mu^{\bullet}(g_{\alpha}f)\) (§1, No.~1, Cor. of Prop. 2). Set \(\nu_{\alpha} = g_{\alpha} \cdot \mu\) ; the condition \(\mu^{\bullet}(gf) < +\infty\) is equivalent to the condition \(\sum_{\alpha \in \mathrm{A}} \nu_{\alpha}(f) < +\infty\) : in other words, \(g\) is locally \(\mu\) -integrable if and only if the family \((\nu_{\alpha})\) is summable. Denoting the sum of this family by \(\nu\) , the preceding calculation yields the equality \(\nu(f) = \mu^{\bullet}(gf)\) , which is equivalent to (8).

COROLLARY. --- Let \((g_{\alpha})\) be a sequence of locally \(\mu\) -integrable numerical functions, such that the sequence of measures \(g_{n} \cdot \mu\) is increasing. In order that this sequence have an upper bound in the ordered vector space \(\mathcal{M}(T)\) of real measures on T, it is necessary and sufficient that the function \(g = \sup g_{n}\) be locally \(\mu\) -integrable; the supremum in \(\mathcal{M}(T)\) of the sequence \((g_{n} \cdot \mu)\) is then the measure \(g \cdot \mu\) .

It suffices to apply Prop. 6 to the functions (positive locally almost everywhere) \(g_n' = g_{n+1} - g_n\) .

PROPOSITION 7. --- Let X be a locally compact space that is countable at infinity, and let \(t \mapsto \lambda_t\) be a \(\mu\) -adequate mapping of T into \(\mathcal{M}_+(X)\) .

Let \(g\) be a positive numerical function defined on \(X\) , locally integrable for the measure \(\nu = \int \lambda_t d\mu(t)\) . The set of \(t \in T\) such that \(g\) is not locally \(\lambda_t\) -integrable is then locally negligible for \(\mu\) , the mapping \(t \mapsto g \cdot \lambda_t\) (defined locally \(\mu\) -almost everywhere) is \(\mu\) -adequate, and

\[
g \cdot \nu = \int (g \cdot \lambda_ {t}) d \mu (t).\tag{9}
\]

Let \((\mathrm{K}_{n})_{n\in\mathbf{N}}\) be an increasing sequence of compact subsets of X whose interiors cover X; if \(\eta\) is any positive measure on X, to say that g is locally \(\eta\) -integrable is equivalent to saying that \(g\varphi_{\mathrm{K}_{n}}\) is \(\eta\) -integrable for every n.~Now let \(\mathrm{H}_{n}\) be the set of \(t\in\mathrm{T}\) such that \(g\varphi_{\mathrm{K}_{n}}\) is not \(\lambda_{t}\) -integrable, and let \(\mathrm{H}=\bigcup_{n}\mathrm{H}_{n}\) ; since \(\mathrm{H}_{n}\) is locally \(\mu\) -negligible for all n (§3, No.~3, Th. 1), the same is true of H, which establishes the first assertion of the statement. Replacing \(\lambda_{t}\) by 0 for t in H (which does not change the measure \(\nu\) ), we can suppose that g is locally \(\lambda_{t}\) -integrable for every \(t\in\mathrm{T}\) . For every \(\nu\) -measurable positive function h defined on X, we have, by Prop. 3 and by Prop. 5 of §3, No.~2,

\[
\int^ {\bullet} h d (g \cdot \nu) = \int^ {\bullet} (g h) d \nu = \int^ {\bullet} d \mu (t) \int^ {\bullet} (g h) d \lambda_ {t} = \int^ {\bullet} d \mu (t) \int^ {\bullet} h d (g \cdot \lambda_ {t}).
\]

This formula and Prop. 5 of §3, No.~2 first show (on taking \(h \in \mathcal{K}_{+}(\mathbf{X})\) ) that the mapping \(t \mapsto g \cdot \lambda_{t}\) is scalarly essentially \(\mu\) -integrable, and that its integral is \(g \cdot \nu\) ; in other words, the relation (9) holds. Next, let us replace \(\mu\) by a positive measure \(\mu' \leqslant \mu\) , and let us take for \(h\) a positive lower semi-continuous function: it follows at once from these relations that \(t \mapsto g \cdot \lambda_{t}\) is \(\mu\) -adequate (§3, No.~1, Def. 1).

PROPOSITION 8. --- Let \(\theta\) be a complex measure on T, \(g_{1}\) a locally \(\theta\) -integrable complex function, \(\theta_{1}\) the measure \(g_{1} \cdot \theta\) . For a complex function \(g_{2}\) defined on T to be locally \(\theta_{1}\) -integrable, it is necessary and sufficient that \(g_{2} g_{1}\) be locally \(\theta\) -integrable, in which case

\[
g _ {2} \cdot \theta_ {1} = g _ {2} \cdot (g _ {1} \cdot \theta) = (g _ {2} g _ {1}) \cdot \theta\tag{10}
\]

(`associativity formula').

By the corollary of Prop. 4, to say that \(g_{2}\) is \(\theta_{1}\) -measurable is equivalent to saying that \(g_{2}g_{1}\) is \(\theta\) -measurable. Let us suppose this condition to be satisfied. For every function \(f \in \mathcal{K}_{+}(T)\) we have, by Propositions 2 and 3,

\[
\int^ {\bullet} | g _ {2} | f d | \theta_ {1} | = \int^ {\bullet} | g _ {2} | f | g _ {1} | d | \theta | = \int^ {\bullet} | g _ {2} g _ {1} | f d | \theta |.
\]

To say that \(g_{2}\) is locally \(\theta_{1}\) -integrable is therefore equivalent to saying that \(g_{2}g_{1}\) is locally \(\theta\) -integrable. Assuming this condition to be satisfied, by Th. 1 we have

\[
\int f d \left(g _ {2} \cdot \theta_ {1}\right) = \int f g _ {2} d \theta_ {1} = \int f g _ {2} g _ {1} d \theta = \int f d \left(g _ {2} g _ {1} \cdot \theta\right),
\]

a formula equivalent to (10).
